\section*{Bab \ref{ch:g12:randvar} --- Peubah Acak dan Distribusi Binomial}

\begin{solution}{exo:g12:randvar:1}
$\E(X) = -0.3 + 0 + 0.8 + 0.5 = 1$.
$\E(X^2) = 0.3 \times 1 + 0 + 0.4\times4 + 0.1\times25 = 4.4$, sehingga menurut
König--Huygens $\V(X) = 4.4 - 1 = 3.4$ dan
$\sigma(X) = \sqrt{3.4} \approx 1.84$.
\end{solution}

\begin{solution}{exo:g12:randvar:2}
\emph{1.} Ada $10$ pertanyaan yang merupakan \omterm{def:g12:randvar:bernoulli}{percobaan Bernoulli} \omterm{def:g12:condprob:indep}{saling bebas} dengan
$p = \frac14$: jadi $X \sim \mathcal B\left(10, \frac14\right)$,
$\E(X) = 2.5$, dan
$\sigma(X) = \sqrt{10 \times \frac14 \times \frac34} = \sqrt{1.875}
\approx 1.37$.

\emph{2.} $\P(X = 0) = \left(\frac34\right)^{10} \approx 0.056$;
\[
\P(X = 5) = \binom{10}{5}\left(\frac14\right)^5\left(\frac34\right)^5
\approx 0.058;
\quad
\P(X \geq 1) = 1 - \left(\tfrac34\right)^{10} \approx 0.944 .
\]
\end{solution}

\begin{solution}{exo:g12:randvar:3}
Percobaan yang \omterm{def:g12:condprob:indep}{saling bebas} dan sejenis: $X \sim \mathcal B(400,\ 0.05)$, sehingga
$\E(X) = 20$ klaim dan
\[
\sigma(X) = \sqrt{400 \times 0.05 \times 0.95} = \sqrt{19} \approx 4.4 .
\]
\end{solution}

\begin{solution}{exo:g12:randvar:4}
Pembayaran $P$ yang diterima memenuhi
$\E(P) = \frac{5 + 6}{6} = \frac{11}{6}$, sehingga harga yang adil adalah
$c = \frac{11}{6} \approx 1.83$ euro. Keuntungan yang adil $G = P - \frac{11}6$
bernilai $-\frac{11}{6}, \frac{19}{6}, \frac{25}{6}$ dengan peluang
$\frac46, \frac16, \frac16$:
\[
\V(G) = \E(G^2)
= \frac{4 \times 121 + 361 + 625}{6 \times 36} = \frac{1470}{216}
= \frac{245}{36} \approx 6.8,
\qquad \sigma(G) \approx 2.6 .
\]
Permainan yang adil berkeuntungan \emph{rata-rata} nol tetapi hasilnya tetap berayun:
jadi adil bukanlah tanpa risiko.
\end{solution}

\begin{solution}{exo:g12:randvar:5}
$X \sim \mathcal B(8,\ 0.7)$.
\[
\P(X = 6) = \binom86 (0.7)^6(0.3)^2 \approx 0.296 ;
\]
\[
\P(X \geq 6) = \P(6) + \P(7) + \P(8)
\approx 0.296 + 8(0.7)^7(0.3) + (0.7)^8
\approx 0.296 + 0.198 + 0.058 = 0.552 ;
\]
$\P(X \geq 1) = 1 - (0.3)^8 \approx 0.99993$. Modusnya:
$(n+1)p = 6.3$, sehingga nilai yang paling mungkin adalah $k^* = 6$
(\cref{exo:g12:randvar:7}).
\end{solution}

\begin{solution}{exo:g12:randvar:6}
Cacah penumpang yang datang adalah $X \sim \mathcal B(104,\ 0.9)$; dan
penerbangannya kelebihan pesanan ketika $X \geq 101$:
\[
\P(X \geq 101) = \sum_{k=101}^{104}\binom{104}{k}(0.9)^k(0.1)^{104-k}
\approx 0.006 .
\]
Menjual $4\%$ tiket lebih banyak daripada kursinya menimbulkan insiden hanya pada kira-kira
$0.6\%$ penerbangannya --- itulah ekonomi di balik kelebihan pesanan itu.
\end{solution}

\begin{solution}{exo:g12:randvar:7}
\[
\frac{\P(X = k+1)}{\P(X = k)}
= \frac{\binom{n}{k+1}}{\binom nk}\cdot\frac{p}{1-p}
= \frac{n - k}{k + 1}\cdot\frac{p}{1-p},
\]
dengan memakai $\binom{n}{k+1} = \binom nk \frac{n-k}{k+1}$. Nisbah ini $\geq 1$
jika dan hanya jika $(n-k)p \geq (k+1)(1-p)$, yaitu $np - k p \geq k - kp + 1 - p$, yaitu
$k \leq (n+1)p - 1$. Jadi peluangnya naik secara tegas selama
$k + 1 \leq (n+1)p$ lalu turun sesudahnya: \omterm{def:g10:functions:extrema}{maksimumnya} tercapai di
$k^* = \floor{(n+1)p}$ (dan dibagi dengan $k^* - 1$ ketika $(n+1)p$ \omterm{def:g10:numbers:sets}{bilangan
bulat}).
\end{solution}

\begin{solution}{exo:g12:randvar:8}
\emph{1.} Gambar pertama pada lemparan $k$ berarti $k-1$ angka lalu gambar, jadi
peluangnya $\left(\frac12\right)^{k-1}\cdot\frac12 = 2^{-k}$, untuk
$1 \leq k \leq 10$. Peluang menang seluruhnya adalah
$\sum_{k=1}^{10} 2^{-k} = 1 - 2^{-10} = \frac{1023}{1024}$ (permainannya
kalah hanya pada sepuluh angka berturut-turut).

\emph{2.} Harapan bayarannya:
\[
\sum_{k=1}^{10} 2^k \cdot 2^{-k} = \sum_{k=1}^{10} 1 = 10 \text{ euro}.
\]
Tanpa batasnya, jumlah $\sum_{k\geq1} 1$ divergen: jadi harapan bayarannya
tak berhingga, walaupun permainannya hampir selalu membayar dengan jumlah kecil --- itulah
\emph{paradoks Saint Petersburg} yang masyhur, yang menunjukkan bahwa harapan saja tidak
mengukur nilai sebuah permainan.
\end{solution}

\begin{solution}{pb:g12:randvar:1}
\textbf{1.} $X \sim \mathcal B\!\left(3, \frac16\right)$:
$\P(X = 0) = \frac{125}{216}$,
$\P(X = 1) = \frac{75}{216}$,
$\P(X = 2) = \frac{15}{216}$,
$\P(X = 3) = \frac{1}{216}$.

\textbf{2.} $\E(X) = 3 \times \frac16 = \frac12$;
$V(X) = 3 \times \frac16 \times \frac56 = \frac{5}{12}$.

\textbf{3.} Harapan bayarannya $2\,\E(X) = 1$ euro terhadap taruhan
$1$ euro: jadi keuntungannya $0$ --- tepat adil (dan itu langka).

\textbf{4.} Tanpa pengembalian pengambilannya bergantungan (sebab
peluang kartu kedua bergantung pada kartu pertama): jadi bukan binomial ---
kotak \omterm{def:g12:condprob:indep}{kebebasan} pada daftar periksanya gagal. Dengan pengembalian:
$n = 5$ yang tetap, $p = \frac14$ yang tetap, pengambilan yang \omterm{def:g12:condprob:indep}{saling bebas}, jadi
binomial.

\textbf{5.} $\E(X) = 6$, $\sigma(X) = \sqrt{4.2} \approx
2.05$. Lalu $\E(S) = 5 \times 6 - 10 = 20$ dan
$\sigma(S) = 5\sigma(X) \approx 10.25$.

\textbf{6.} Ada $n = 105$ percobaan yang tetap (yaitu tiketnya), masing-masing berupa
datang atau tidak dengan $p = 0.9$ yang sama, dan dianggap \omterm{def:g12:condprob:indep}{saling bebas}: jadi
binomial. Pengakuannya: keluarga ketinggalan pesawat
\emph{bersama-sama} dan badai mengosongkan seluruh pesawat --- jadi \omterm{def:g12:condprob:indep}{kebebasannya}
adalah lompatan iman modelnya, dan ketidakhadiran yang berkorelasi membuat
ekornya lebih gemuk daripada yang dijanjikan \omterm{def:g12:randvar:binomial}{distribusi binomial}.

\textbf{7.} $\E(X) = 94.5$;
$\sigma(X) = \sqrt{105 \times 0.9 \times 0.1} =
\sqrt{9.45} \approx 3.07$.

\textbf{8.} $\P(X \geq 101) = \sum_{k=101}^{105}
\binom{105}{k} 0.9^k\, 0.1^{105 - k}$.

\textbf{9.} Nilainya $\approx 0.0167$: jadi kira-kira satu penerbangan dari enam puluh
menggusur seseorang.

\textbf{10.} $\E(\text{tergusur}) \approx 0.023$ penumpang tiap
penerbangan, sehingga harapan ganti ruginya $\approx 19$ euro --- terhadap
pemasukan tambahan $1\,000$ euro. Kelebihan pesanan sebanyak lima itu
menguntungkan luar biasa; karena itulah setiap maskapai melakukannya.

\textbf{11.} Dari tabel peluang ekornya:
$n = 105$: $1.7\,\%$; $n = 106$: $4.0\,\%$; $n = 107$:
$8.1\,\%$. Jadi penjualan terbesar yang taat aturan adalah $n = 106$ tiket.

\textbf{12.} $z = \frac{100.5 - 94.5}{3.07} \approx 1.95$:
jadi ambangnya duduk dua \omterm{def:g11:stat:variance}{simpangan baku} di atas rata-ratanya,
dan kaidah kasar kurva lonceng (``ekor atas $2\sigma$
$\approx 2.5\,\%$'') mendarat dekat dengan nilai eksaknya, $1.7\,\%$ ---
kurva mulus yang membayangi \omterm{def:g12:randvar:binomial}{distribusi binomial} itu adalah tokoh utama
bab berikutnya.

\textbf{13.} $0.98^{20} + 20 \times 0.02 \times 0.98^{19}
\approx 0.94$: jadi lot yang jujur lolos $94\,\%$ waktunya.

\textbf{14.} $0.9^{20} + 20 \times 0.1 \times 0.9^{19}
\approx 0.39$: jadi lot yang buruk masih menyelinap lolos $39\,\%$
waktunya.

\textbf{15.} Risiko produsennya: lot yang baik ditolak
($\approx 6\,\%$); risiko konsumennya: lot yang buruk diterima
($\approx 39\,\%$). Contoh yang lebih besar (dengan ambang
penerimaan yang sebanding) menyusutkan keduanya --- dengan harga pemeriksaan
yang lebih banyak: jadi mutu punya garis anggarannya.

\textbf{16.} $V\!\left(\frac Xn\right) = \frac{V(X)}{n^2} =
\frac{np(1-p)}{n^2} = \frac{p(1-p)}{n}$, jadi \omterm{def:g11:stat:variance}{simpangan bakunya}
$\sqrt{\frac{p(1-p)}{n}}$. Lipatempatkan contohnya, separuhkan
deraunya: itulah hukum $\frac{1}{\sqrt n}$ pada selang ayunan
kelas 10, yang kini terturunkan.

\textbf{17.} Batasnya mengekang bayarannya pada $2^{10} = 1024$
euro, sehingga masing-masing dari sepuluh putarannya menyumbang tepat $1$ euro
harapan: $\E = 10$. Harapan tak berhingga pada permainan tanpa batas itu
seluruhnya berasal dari bayaran yang luar biasa jarang dan
luar biasa besar; membatasinya berarti mengakui bahwa tak ada bank
yang membayar $2^{50}$ euro. Harga tiket yang adil bagi permainan berbatas ini:
$10$ euro --- dan tiba-tiba tak ada lagi yang terparadoks.

\textbf{18.} Harapan uang hadiah tiap tiketnya:
$\frac{100 + 2 \times 50}{200} = 1$ euro terhadap tiket
$2$ euro: jadi marjinnya $50\,\%$ --- delapan belas kali marjin rolet yang
$2.7\,\%$. Undiannya bertahan sebab pemainnya
sadar sedang menyumbang: ``kerugiannya'' justru itulah maksudnya.

\textbf{19.} Harapan klaimnya: $1\,000 \times 0.01 \times
10\,000 = 100\,000$; preminya $130\,000$, jadi harapan labanya
$30\,000$ euro. \omterm{def:g11:stat:variance}{Simpangan baku} total klaimnya:
$10\,000 \times \sqrt{1\,000 \times 0.01 \times 0.99}
\approx 31\,500$ euro --- jadi satu tahun buruk yang biasa saja melahap
seluruh harapan labanya. Karena itulah ada cadangan modal, penanggungan ulang,
dan portofolio yang jauh lebih besar daripada seribu polis: penanggung asuransi
hidup dari hukum bilangan besar dan memegang cadangan terhadap
kelambatannya.

\textbf{20.} Daftar periksa dahulu --- $n$ yang tetap, $p$ yang sama,
\omterm{def:g12:condprob:indep}{kebebasan} yang \emph{diklaim} --- atau tidak ada binomial sama sekali. Lalu
berlayarlah dengan $\E \pm \sigma$: rata-ratanya menempatkan, simpangannya memperingatkan.
Lalu hargailah ekornya: penggusuran, lot yang buruk, dan tahun yang membangkrutkan semuanya
hidup di luar $2\sigma$. Peringatannya: \omterm{def:g12:condprob:indep}{kebebasan} itu janji sang pemodel,
bukan janji dunia; dan harapan justru mengabaikan
hal yang menghancurkanmu. Halaman yang hilang dari buku aturannya --- seberapa cepat
frekuensinya mengendap, dan bentuk apa yang diambil ayunannya ---
adalah kedua bab berikutnya.

\end{solution}
